% ============================================================
%  sections/conclusion.tex
% ============================================================

\section{Conclusion}

We presented \modelname{}, a modality-aware adaptive token routing
framework that trains lightweight sparsity masks end-to-end using a
differentiable top-\(k\) relaxation. Our experiments on the new
\dataset{} benchmark demonstrate \textbf{3.1\(\times\)} FLOPs and
\textbf{2.4\(\times\)} latency reductions over the full-attention
\baseline{} baseline at negligible accuracy cost on four
vision--language tasks. Ablations confirm that modality-specific masks
and the residual injector are both essential, and that masks
transfer reliably across ViT scales without retraining.

\textbf{Future work} includes extending adaptive routing to autoregressive
decoding, exploring mask sharing across layers to reduce parameter
overhead, and integrating with quantisation for further efficiency gains.
We will release code, model checkpoints, and the \dataset{} evaluation
harness upon acceptance.

\vspace{8pt}
\noindent
\begin{tcolorbox}[methodbox, title={Poster Highlights}]
  \small
  \begin{itemize}[leftmargin=1em, itemsep=3pt, topsep=2pt, label=\textbullet]
    \item \textbf{3.1\(\times\)} FLOPs reduction vs.\ ViT-L/14 baseline
    \item \textbf{2.4\(\times\)} wall-clock latency on A100-80G
    \item \(\leq\)0.5\% accuracy drop across all 12 \dataset{} tasks
    \item Zero-shot mask transfer across ViT-B, ViT-L, ViT-H
    \item Compatible with existing checkpoints --- no backbone surgery
  \end{itemize}
\end{tcolorbox}

\section*{Acknowledgements}

The authors thank the Vertex Institute HPC team for compute access,
and the anonymous reviewers for their constructive feedback.
This work was partially supported by Nimbus Systems under Research
Collaboration Agreement NMS-2024-117.
